\scanpage{188}
\mspage{137}
\unsure{et} il s'agit donc de comparer les deux intégrales
\diagram{\begin{tikzcd}[ampersand replacement=\&,row sep=tiny,column sep=small]
\&X_1\&\&X_1^0\\
\mathscr X=\int X_0\arrow[ru,-,"\alpha_1"]\arrow[rd,"\alpha_2", "\deleted{\unsure{\alpha}}"']\&{,}\&\widetilde{\mathscr X}^0=\int X_0^0\arrow[ru,"\alpha_1^0"]\arrow[rd,"\alpha_2^0"]\&{,}\\
\&X_2\&\&X_2^0
\end{tikzcd}}
\unsure{notamment} comme objets de $\mathrm{HOT}_W(\Phi)$.

Pour cela, je reprends la construction \unsure{d'un|du} foncteur
\[
X\unsure{\longmapsto|\longleftarrow}\Phi(X):\quad\mathit{Cat}\longrightarrow\mathit{Cat}
\]
de I, p.~72\deleted{\unsure{,}}-75, mais je préfère ici noter ce foncteur $\Sigma(X)$ (car $\Phi$ désigne déjà une \deleted{petite} catégorie
\diagram{\begin{tikzcd}[ampersand replacement=\&,row sep=tiny,column sep=tiny]
\&1\\0\arrow[ru,-]\arrow[rd,-]\&\\\&2
\end{tikzcd}}
). On utilisera \unsure{l'exist. de|l'exist.|l'existence} homs fonctoriels en $X$ :
\diagram{\begin{tikzcd}[ampersand replacement=\&,row sep=small,column sep=small]
\&\Sigma(X)\arrow[dl,"p_X"']\arrow[dr,"q_X"]\&\\
X\&\&X^0
\end{tikzcd}}
qui sont $W$-asphériques, donc $\in W$. De plus, $\Sigma(X^0)=\Sigma(X)$, $p_{X^0}=q_X$, $q_{X^0}=p_X$.

Ainsi, \deleted{\illegible{1}} si
\diagram{\begin{tikzcd}[ampersand replacement=\&,row sep=tiny,column sep=small]
\& X_1\&\\
X_{\ast}=\Bigl(\,X_0\arrow[ru,"\alpha_1"]\arrow[rd,"\alpha_2"]\&\Bigr)\&\in\mathrm{Hom}(\Phi^0,\mathit{Cat})\\
\& X_2\&
\end{tikzcd}}
\diagram{\begin{tikzcd}[ampersand replacement=\&,row sep=tiny,column sep=small]
\&X_1^0\&\\
X_{\ast}^{0}=\Bigl(\,X_0^0\arrow[ru,-,"\alpha_1^0"]\arrow[rd,"\alpha_2^0"]\&\Bigr)\&\in\quad \unsure{\mathrm{id}|\mathrm{id}^{\mathrm{em}}}\rlap{,}\\
\&X_2^0\&
\end{tikzcd}}
% archive sheet number: 188
